% =====================================================================
%  Un refactoring pour le portage, une transmission
% =====================================================================
\begin{frame}{Un refactoring pour le portage, une transmission}
  \vspace{0.2cm}
  \begin{columns}[c,onlytextwidth]
    \column{0.50\textwidth}
      \centering
      \resizebox{\textwidth}{!}{%
      \begin{tikzpicture}[font=\sffamily]
        \node[flowhi,text width=3.2cm,minimum height=1.6cm] (core)
          {\flowtxthi{\emph{core}}{le socle commun}};
        \foreach \i/\lab [count=\n from 0] in
          {1/{fork 1},2/{fork 2},3/{fork 3},4/{fork 4},5/{fork 5},6/{fork 6},7/{prépayé}}{
          \node[draw=accent,line width=1pt,fill=white,rounded corners=4pt,
                minimum width=2.6cm,minimum height=0.78cm,font=\bfseries\fontsize{11}{13}\selectfont,
                text=ink] (s\i) at (6.4,4.2-1.05*\n) {\emph{spec} · \lab};
          \draw[flowarr,line width=1.2pt] (core.east) to[out=0,in=180] (s\i.west);
        }
      \end{tikzpicture}}
    \column{0.46\textwidth}
      \card{2.3cm}{Une refonte \emph{core} / \emph{spec}}{Un socle commun, puis une
        \emph{spec} par fork qui n'isole que les différences métier.}
      \vspace{0.28cm}
      \cardhi{2.3cm}{Une transmission assumée}{Mené en \textbf{binôme}, pour
        transmettre la connaissance du framework avant mon départ.}
  \end{columns}
\end{frame}
